\SourceLeaf{129}
bloss oberflächliche Aufmerksamkeit zu widmen, um besonders den praktischen Werth der davon zu hoffenden Vortheile würdigen zu können.

Der Vf. hatte seine Schrift dem französischen Nationalinstitute zur Beurtheilung vorgelegt; der Bericht, welchen \textsc{Delambre} darüber abgestattet und welchen das Institut gebilligt hat, ist der Schrift selbst beigefügt. \textsc{Leonelli} ist damit nicht ganz zufrieden gewesen, und äussert in seinen gleichfalls angehängten Gegenbemerkungen seine Empfindlichkeit einigemal mit vieler Lebhaftigkeit. Wir werden nicht umhin können, auch \textsc{Delambres} Bericht und \textsc{Leonellis} Beantwortung mit zu berühren, da beide nun einen Theil des Buchs selbst ausmachen.

Bei der in der \emph{ersten} Abtheilung vorgetragenen Methode, den Logarithmen jeder vorgegebenen Grösse $A$ zu berechnen, liegt die Hauptidee zum Grunde, dass diese Grösse in ein Product von der Form
\[
10^{\mu} a \left(1+\frac{1}{10} b\right)\left(1+\frac{1}{100} c\right)\left(1+\frac{1}{1000} d\right)\ldots
\]
verwandelt werde, so dass $a$, $b$, $c$, $d\ldots$ einfache ganze Zahlen bedeuten. Man sieht erstlich, dass die Logarithmen dieser einzelnen Factoren eine abnehmende Reihe bilden werden, von denen man sich mit einer gewissen grössern oder kleinern Zahl begnügen kann, je nachdem man den Logarithmen des Products mit mehr oder weniger Decimalstellen verlangt; zweitens dass man alle jene einzelnen Logarithmen sogleich in einer ein für allemal berechneten Tafel vorräthig haben kann. Eine solche Tafel liefert \textsc{Leonelli} hier zuerst für die Briggischen Logarithmen auf 20 Decimalstellen. Diese enthält zuerst die Logarithmen der ganzen Zahlen von 2 bis 9; sodann die Logarithmen für alle $1+\frac{1}{10} b$, d.\,i. für 1,1; 1,2; 1,3 bis 1,9; nachher für alle $1+\frac{1}{100} c$ u.\,s.\,w. Ausser den 8 ganzen Zahlen kommen also, für jede Ordnung der gebrochenen, 9 Logarithmen; weiter als bis zur 11$^{\mathrm{ten}}$ Ordnung brauchte aber die Tafel nicht ausgedehnt zu werden, da bei der zwölften die Logarithmen mit den Logarithmen der eilften einerlei bedeutende Ziffern haben, und so bei den folgenden. Auf diese Weise umfasst die Tafel nur 107 Logarithmen, welche hinreichen, um den Logarithmen jedes Products der obigen Form auf 18 oder 19 Ziffern durch blosse Addition zu berechnen; die 20$^{\mathrm{ste}}$ Stelle bleibt dabei offenbar immer schwankend. Für die hyperbolischen Logarithmen hat \textsc{Leonelli} eine ähnliche Tafel von demselben Umfange beigefügt.

